\section{Scene IR: The Universal Render Contract}
\label{sec:scene-ir}

The Scene IR is the central abstraction of the kernel---the single handoff point between any grammar's layout engine and every rendering backend. It is a backend-agnostic, byte-deterministic description of every drawing primitive, visual treatment, and effect request.

\subsection{Design Rationale}

\paragraph{Why a Scene IR?}
A Scene IR (sometimes called a ``Render IR'' or ``Display List'') decouples the concerns of layout from the concerns of output format. This separation enables:

\begin{itemize}
  \item \textbf{Grammar-agnosticism}: Timeline, flow, and graph grammars all compile to the same Scene IR. The backends need not know which grammar produced the scene.
  
  \item \textbf{Backend-agnosticism}: SVG, PNG, PDF, and future backends consume the same input. Adding a new backend requires no changes to any grammar.
  
  \item \textbf{Determinism verification}: The Scene IR is hashable; golden-testing can verify that the same IR+theme produces byte-identical scenes.
  
  \item \textbf{Effect negotiation}: Effects (glow, drop-shadow, animation) are \emph{requested} in the Scene IR with fallback policies. Backends fulfill or degrade gracefully.
\end{itemize}

\subsection{Scene Model}
\label{sec:scene-model}

The Scene consists of four top-level structures:

\begin{Verbatim}[frame=single,fontsize=\small]
Scene {
  canvas:    { width, height, background }
  elements:  [ DrawingPrimitive ]   -- ordered; painting order preserved
  effects:   { effect_id -> EffectDefinition }
  meta:      { theme_id, fidelity_tier, scene_hash(SHA-256), version }
}
\end{Verbatim}

\subsubsection{Canvas Descriptor}

\begin{Verbatim}[frame=single,fontsize=\small]
Canvas {
  width:      number       -- logical pixels
  height:     number       -- computed from content
  background: Background   -- solid color, gradient, or pattern
}

Background := one of:
  SolidColor     { color: "#RRGGBB" | "#RRGGBBAA" }
  LinearGradient { angle_deg, stops: [(offset, color)...] }
  RadialGradient { cx, cy, radius, stops: [(offset, color)...] }
\end{Verbatim}

\subsubsection{Drawing Primitives}

The Scene IR defines a minimal set of drawing primitives---the ``assembly language'' of visual output:

\begin{Verbatim}[frame=single,fontsize=\small]
DrawingPrimitive := one of:

  Rect {
    id, x, y, width, height,
    fill, stroke, stroke_width, corner_radius,
    opacity, clip_id?, effect_refs[]
  }

  Circle {
    id, cx, cy, radius,
    fill, stroke, stroke_width, opacity, effect_refs[]
  }

  Line {
    id, x1, y1, x2, y2,
    stroke, stroke_width, stroke_style,   -- solid | dashed | dotted
    opacity, effect_refs[], animation_ref?
  }

  Path {
    id, d_commands: [(op, args)...],      -- SVG path mini-language
    fill, stroke, stroke_width, opacity, effect_refs[]
  }

  Text {
    id, x, y, content,
    font_family, font_size, font_weight, color,
    anchor: 'start' | 'middle' | 'end',
    clip_id?, effect_refs[]
  }

  MultiText {
    id, x, y, lines: [TextSpan...],
    line_height, clip_id?
  }

  Image {
    id, x, y, width, height,
    data_uri: string                      -- embedded raster or SVG
  }

  Group {
    id, children: [DrawingPrimitive...],
    clip_rect?, transform?, opacity?
  }
\end{Verbatim}

\paragraph{No semantic primitives.}
The Scene IR contains \emph{no} semantic types (Activity, Milestone, Track, Node, Edge). By the time content reaches the Scene, it has been reduced to pure geometry. This is the key to grammar-agnosticism.

\subsubsection{Effect Definitions}

Effects are visual treatments that may not be uniformly supported across backends. The Scene IR \emph{requests} effects and specifies fallback policies:

\begin{Verbatim}[frame=single,fontsize=\small]
EffectDefinition := one of:

  Glow {
    color, radius_px, intensity,
    fallback_policy: 'approximate' | 'omit' | 'embed-raster' | 'error'
  }

  DropShadow {
    offset_x, offset_y, blur_radius, color, opacity,
    fallback_policy
  }

  GradientFade {
    direction: 'left' | 'right' | 'top' | 'bottom',
    color_stops: [(offset, color)...],
    fade_px,
    fallback_policy
  }

  NoiseTexture {
    seed, scale, turbulence_type, opacity, color,
    fallback_policy
  }

  Bloom {
    radius_px, threshold, intensity,
    fallback_policy
  }
\end{Verbatim}

\paragraph{Fallback policies.}
Each effect carries a \texttt{fallback\_policy} instructing backends what to do when they cannot render the effect natively:

\begin{center}
\small
\begin{tabular}{ll}
\toprule
\textbf{Policy} & \textbf{Meaning} \\
\midrule
\texttt{approximate} & Use the closest available native mechanism \\
\texttt{omit}        & Silently omit the effect; geometry unchanged \\
\texttt{embed-raster}& Rasterize the affected layer; embed as Image \\
\texttt{error}       & Abort with a descriptive diagnostic \\
\bottomrule
\end{tabular}
\end{center}

\subsection{Animation Hints}
\label{sec:scene-animation}

Animation is modeled as \emph{optional, declarative hints} attached to primitives. Backends that support animation (SVG, HTML) honor the hints; backends that produce static output (PNG, PDF) render the resting frame.

\begin{Verbatim}[frame=single,fontsize=\small]
AnimationHint := one of:

  FlowingDashes {
    target_id,                            -- Line or Path id
    dash_length, gap_length,
    speed_px_per_sec,
    direction: 'forward' | 'reverse'
  }

  DrawOn {
    target_id,                            -- Path id
    duration_ms,
    easing: 'linear' | 'ease-in-out' | 'ease-out'
  }

  Pulse {
    target_id,
    property: 'opacity' | 'scale' | 'fill',
    from, to, duration_ms, repeat: number | 'infinite'
  }

  DotAlongPath {
    target_id,                            -- Path id
    dot_radius, dot_fill,
    duration_ms, repeat
  }
\end{Verbatim}

\paragraph{Determinism is preserved.}
An animated SVG is still byte-identical markup---the animation is declarative SMIL or CSS, not frame-by-frame. Golden tests can compare the SVG text. A rasterized frame is a static snapshot of the resting state.

\subsection{Scene Determinism Contract}
\label{sec:scene-determinism}

The Scene is byte-deterministic: two executions of any layout engine on identical IR and theme produce bitwise-identical Scene records.

\paragraph{The \texttt{scene\_hash}.}
The Scene's metadata includes a SHA-256 hash of the serialized element list and effect registry. This hash is reproducible across runs and is used for:

\begin{itemize}
  \item Golden-image testing (compare hashes before comparing pixels)
  \item Cache invalidation (skip re-render if hash matches)
  \item Audit trails (record which scene produced which artifact)
\end{itemize}

\paragraph{Layered determinism.}
Determinism applies at three levels:

\begin{enumerate}
  \item \textbf{Scene geometry---always byte-deterministic.} All coordinates, dimensions, and element ordering are bitwise-reproducible.
  
  \item \textbf{Per-backend output---deterministic given pinned backend version.} The SVG backend with version X produces identical output from identical scenes. Effect seeds for procedural effects (noise, cloud) derive from \texttt{scene\_hash + effect\_id}.
  
  \item \textbf{Cross-backend pixel identity---explicitly not promised.} The SVG and Skia backends produce visually different outputs (Skia renders art effects; SVG approximates or omits). This is correct behavior, not a defect.
\end{enumerate}

\subsection{Scene Serialization}

The Scene IR has a canonical JSON serialization for debugging, testing, and tooling interoperability. The serialization is:

\begin{itemize}
  \item \textbf{Sorted keys} at every level (for diff-friendliness)
  \item \textbf{No whitespace-only differences} (normalized indentation)
  \item \textbf{Floating-point precision} capped at two decimal places
\end{itemize}

This serialization is the basis for the \texttt{scene\_hash} computation.

\subsection{Scene IR as the Integration Point}

The Scene IR is the \emph{only} integration point between the upper layers (grammars, composition) and the lower layers (backends, output formats). This architectural choice has several consequences:

\begin{description}
  \item[Adding a grammar] requires only implementing a layout engine that emits Scene IR. No backend changes are needed.
  
  \item[Adding a backend] requires only consuming Scene IR. No grammar changes are needed.
  
  \item[Testing] can happen at the Scene level: grammar tests verify Scene output; backend tests verify rendering from canonical scenes.
  
  \item[Themes] operate entirely within the Scene---they determine fill colors, stroke widths, and effect parameters, but they do not affect layout (which is already baked into coordinates).
\end{description}
